\section{Limited systems and Lag control}

\subsection{Ctrl+F keywords}

\subsubsection{Saturation, rate limit, actuator limit, Lag, P-Lead-Lag, beta, begraensning}
\[
\text{amplitude limit, rate limit, windup, lag compensator, } \beta,N_i,\tau_i
\]
Use this section when the task includes actuator limits or a P-Lead-Lag controller.

\subsection{Actuator limitations}

\subsubsection{Amplitude saturation}
\[
u_{\mathrm{sat}}(t)=
\begin{cases}
u_{\max}, & u(t)>u_{\max},\\
u(t), & u_{\min}\leq u(t)\leq u_{\max},\\
u_{\min}, & u(t)<u_{\min}
\end{cases}
\]
Saturation is nonlinear and can invalidate a linear design.

\subsubsection{Rate limitation}
\[
|\dot{u}(t)|\leq R_{\max}
\]
Rate limits prevent fast command changes and can reduce effective phase margin.

\subsubsection{Integrator windup recognition}
\[
u_I(t)=\frac{K_P}{\tau_i}\int_0^t e(\xi)\,d\xi
\]
If the actuator saturates while the integral state keeps growing, expect overshoot and slow recovery.

\subsection{Lag controller}

\subsubsection{Lag transfer function}
\[
C_{\mathrm{lag}}(s)=\frac{\tau_i s+1}{\frac{\tau_i}{\beta}s+1},\qquad \beta>1
\]
Lag increases low-frequency gain relative to high-frequency gain without adding an integrator.

\subsubsection{P-Lead-Lag controller}
\[
C(s)=K_P
\frac{\tau_d s+1}{\alpha\tau_d s+1}
\frac{\tau_i s+1}{\frac{\tau_i}{\beta}s+1}
\]
This controller combines phase lead with low-frequency gain shaping.

\subsubsection{Lag phase contribution at crossover}
\[
\phi_{\mathrm{lag}}(\omega_c)=
\arctan\left(\frac{N_i(1-\beta)}{1+\beta N_i^2}\right),
\qquad N_i=\omega_c\tau_i
\]
This formula is directly tied to the recurring P-Lead-Lag design task.

\subsubsection{Lag steady-state error}
\[
e_{\mathrm{ss}}=\frac{1}{1+L(0)}
\]
Lag gives finite low-frequency gain, so zero step error is not automatic.

\subsection{Exam recipe: solve for Lag parameter \(\beta\)}

\textbf{Use when:} the problem gives a P-Lead-Lag controller and asks for \(\beta\) or lag phase.

\textbf{Steps:}
\begin{enumerate}
\item Read or compute \(N_i=\omega_c\tau_i\).
\item Determine the allowed or required lag phase \(\phi_{\mathrm{lag}}\).
\item Insert \(N_i\) and \(\phi_{\mathrm{lag}}\) in the lag phase equation.
\item Solve for \(\beta>1\).
\item Check the resulting low-frequency gain and phase-margin loss.
\end{enumerate}

\textbf{Fast checks:} lag phase is negative; if the result gives \(\beta\leq1\), it is not a lag compensator.

\textbf{Common traps:} treating Lag like PI; Lag has no pole at the origin and does not guarantee zero step error.

\subsection{Exam recipe: recognize limit-induced design problems}

\textbf{Use when:} the task asks why a linear design fails on a real actuator.

\textbf{Steps:}
\begin{enumerate}
\item Compare requested \(u(t)\) with amplitude limits.
\item Compare requested \(\dot{u}(t)\) with rate limits.
\item Check whether integral action continues during saturation.
\item Reduce bandwidth, filter reference, or add anti-windup if limits dominate.
\end{enumerate}

\textbf{Fast checks:} large \(K_P\), derivative action, and fast prefilters tend to increase actuator demand.
